\documentclass[11pt]{article}
\usepackage[a5paper,margin=16mm]{geometry}
\usepackage{fontspec}
\setmainfont{TeX Gyre Termes}
\setsansfont{TeX Gyre Heros}
\usepackage{amsmath,amsthm}
\usepackage{graphicx}
\usepackage[colorlinks=true,linkcolor=blue,citecolor=teal,urlcolor=magenta,
  pdftitle={Cross-reference matrix},pdfauthor={Ratex fixture},bookmarksnumbered=true]{hyperref}
\usepackage[capitalise,noabbrev]{cleveref}
\newtheorem{theorem}{Theorem}[section]
\crefname{theorem}{Theorem}{Theorems}

\begin{document}
\tableofcontents
\section{Setup}\label{sec:setup}
We state \cref{thm:main} in \cref{sec:results}; its proof uses \cref{eq:a,eq:b}.
The summary appears on \cpageref{sec:end} and the table in \cref{tab:data}.
\begin{equation}\label{eq:a} a^{2}+b^{2}=c^{2}\end{equation}
\begin{equation}\label{eq:b} c = \sqrt{a^{2}+b^{2}}\end{equation}
\clearpage
\section{Results}\label{sec:results}
\begin{theorem}\label{thm:main}
\Cref{eq:a} and \cref{eq:b} are equivalent for $c\ge0$.
\end{theorem}
\subsection{Numbers}\label{subsec:numbers}
\begin{table}[ht]
\centering\caption{Counts}\label{tab:data}
\begin{tabular}{lr}\hline x & 1\\ y & 2\\\hline\end{tabular}
\end{table}
\begin{figure}[ht]
\centering\fbox{\rule{0pt}{1.2cm}\rule{3cm}{0pt}}
\caption{A placeholder box}\label{fig:box}
\end{figure}
\Cref{fig:box,tab:data} are floats; \cref{sec:setup,subsec:numbers,sec:results} are sections.
Jump to \nameref{sec:end} (\autoref{sec:end}) or visit \url{https://example.org/ratex}.
\footnote{Footnotes keep hyperlinks: see \hyperref[sec:setup]{the setup}.}
\clearpage
\section{Summary}\label{sec:end}
\Cref{thm:main} on \cpageref{thm:main}; \cref{eq:a} through \cref{eq:b} as a range:
\crefrange{eq:a}{eq:b}.
\end{document}
